\documentclass[10pt,a4paper]{amsart}
\usepackage[margin=19mm]{geometry}
\usepackage{amsmath,amssymb,mathtools}
\usepackage[scr=rsfs,cal=euler]{mathalfa}
\usepackage{xfrac}
\usepackage[unicode,hidelinks,bookmarks=false]{hyperref}
\newcommand{\ie}{i.e.\ }
\newcommand{\cf}{cf.\ }
\newcommand{\resp}{respectively\ }
\newcommand{\Subsection}{\subsection}
\providecommand{\nobreakdash}{\nobreak\hbox{-}}
\setlength{\emergencystretch}{2em}
\multlinegap0pt
\usepackage{bm}
\usepackage{bbm}
\usepackage{dsfont}
\usepackage{xspace}
\usepackage{cancel}
\usepackage{mathtools}
\usepackage{amsthm}
\makeatletter
\@ifundefined{theorem}{\newtheorem{theorem}{Theorem}}{}
\@ifundefined{lemma}{\newtheorem{lemma}[theorem]{Lemma}}{}
\makeatother
\DeclareMathOperator{\KL}{KL}
\newcommand{\Xx}{\mathcal{X}}
\renewcommand{\d}{\mathrm{d}}
\title[Doubly Regularized Entropic Wasserstein Barycenters]{Doubly Regularized Entropic Wasserstein Barycenters}
\author{L\'ena\"ic Chizat}
\date{Source: arXiv:2303.11844v2 (CC BY 4.0)}
\begin{document}
\maketitle

\noindent\emph{Source and licence.} Doubly Regularized Entropic Wasserstein Barycenters. Version arXiv:2303.11844v2 (CC BY 4.0), distributed under the Creative Commons Attribution 4.0 International licence, as recorded in the arXiv licence metadata for this exact version and shown on the abstract page https://arxiv.org/abs/2303.11844v2. Full rights notice: licenses/arxiv-2303.11844-CC-BY-4.0.txt.\par\smallskip

\noindent\textit{Editorial scope.} Counted unit: the Lemma whose source label is \texttt{lem:positive-minimizer} in the article (source0.tex lines 1173--1175, source label \texttt{lem:positive-minimizer}), delivered together with the article's own complete original proof of it (source0.tex lines 1176--1212, one \texttt{proof} environment of 37 lines). No mathematics is written, repaired or re-derived: statement and proof are the article's own text line for line. Required context retained uncounted: none. Every definition, display and label of those lines is retained unchanged. Omitted from this package and not redistributed: 1--1172, 1213--1235. Nothing inside a retained excerpt is omitted or altered: every retained body is delivered exactly as the article's lines are written, so no omission receipt of any kind accompanies this package. Citations: the retained excerpts carry no citation command at all, so no bibliography entry of the article is transcribed and no reference is redistributed. Reference adaptation: none was needed and none was made; every reference inside the delivered unit resolves inside it, so no unresolved reference or citation is produced. Printed slips kept literal: no printed word, symbol, hypothesis or number of the counted statement or of its proof was altered. Layout and class: the article's own class is \texttt{article} with packages including inputenc, geometry, enumitem, mathtools, amsmath, amssymb, graphicx, url, amsthm, subcaption, dsfont, mathrsfs, algorithm2e, biblatex; the wrapper is the lane's pinned \texttt{amsart} editorial wrapper at 10pt on A4, which restates the theorem-like environments the retained text needs and adds the article's own macro definitions that the retained text needs (\texttt{\textbackslash KL}, \texttt{\textbackslash Xx}, \texttt{\textbackslash d}), with extra wrapper packages (bm, bbm, dsfont, xspace, cancel, mathtools). The delivered numbering is the unit's own. Assets: no retained span contains a figure environment, a graphics-inclusion command, a PDF-page inclusion, a table or a TikZ picture; no image file is copied into this package, the package declares no asset dependency and no image file is redistributed with it. Licence: version arXiv:2303.11844v2 of this article is distributed under the Creative Commons Attribution 4.0 International licence, as recorded in the arXiv licence metadata for this exact version and shown on the abstract page https://arxiv.org/abs/2303.11844v2. A full rights notice is delivered at licenses/arxiv-2303.11844-CC-BY-4.0.txt. Characters: every retained excerpt is pure ASCII, so the delivered engine, XeLaTeX with its default Latin Modern text font, reports no missing character.
\begin{lemma}\label{lem:positive-minimizer}
Any minimizer $\mu^*$ of $F$ is absolutely continuous with respect to $\pi$ and its density $\rho^*=\frac{\d \mu^*}{\d \pi}$ must satisfy $\rho^*>0$ and $\log \rho^*\in L^1(\pi)$.
\end{lemma}

\begin{proof}
For $\epsilon\in [0,1]$, let us define $\rho_\epsilon = (1-\epsilon)\rho^*+\epsilon$, $\mu_\epsilon=\rho_\epsilon \pi$ and compare $F(\mu^*)$ to $F(\mu_\epsilon)$. By optimality of $\mu^*$ and convexity of $G$, we may write
$$
\KL(\mu^*|\pi)- \KL(\mu_\epsilon|\pi) \leq G(\mu_\epsilon) - G(\mu^*)\leq (1-\epsilon) G(\mu^*) +\epsilon G(\pi) - G(\mu^*) = \epsilon (G(\pi)-G(\mu^*)).
$$
We get 
$$
\int ( f(\rho^*)-f(\rho_\epsilon))\d\pi \leq C\epsilon
$$
where $f(t)=t\log(t)-t+1$ (and $f(0)=1$ by convention). Write 
\begin{align*}
A = \{x\in \Xx\;;\; \rho^*(x)>0\}, && B = \{x\in \Xx\;;\; \rho^*(x)=0\}.
\end{align*}
Since $f$ is convex we have for $x\in A$
$$
f(\rho^*(x))-f(\rho_\epsilon(x)) \geq ((\rho^*(x)-\rho_\epsilon(x))f'(\rho_\epsilon(x)) = \epsilon (\rho^*(x)-1)\log \rho_\epsilon(x).
$$
For $x\in B$, we simply write $f(\rho^*(x))-f(\rho_\epsilon(x)) =f(0)-f(\epsilon)=-\epsilon \log \epsilon -\epsilon$. This allows to write
$$
-\epsilon (\log \epsilon +1)\pi(B) +\epsilon  \int_{A} (\rho^*(x)-1)\log \rho_\epsilon(x) \d\pi(x) \leq C\epsilon
$$
and, dividing by $\epsilon$,
\begin{align}\label{eq:proof-positive}
-(\log \epsilon +1)\pi(B) +  \int_{A} (\rho^*(x)-1)\log \rho_\epsilon(x) \d\pi(x) \leq C.
\end{align}
Note that we always have $(\rho^*(x)-1)\log \rho_\epsilon(x)\geq 0$ (just distinguish between the case $\rho^*(x)\geq 1$ and $\rho^*(x)\leq 1$). Thus, we may write
$$
-(\log \epsilon +1)\pi(B) \leq C.
$$
Letting $\epsilon\to 0$ proves $\pi(B)=0$, hence $\rho^*(x)>0$, $\pi$-almost everywhere.

We now come back to Eq.~\eqref{eq:proof-positive} which is an upper-bound on the nonnegative functions $(\rho^*(x)-1)\log \rho_\epsilon(x)$. By Fatou's lemma we have as $\epsilon\to 0$
$$
\int_{\Xx} (\rho^*(x)-1)\log \rho^*(x) \d\pi(x) \leq C.
$$
Since this is the integral of a nonnegative function, it follows that $(\rho^*-1)\log \rho^*$ is in $L^1(\pi)$. Since we already know that $\rho^*\log \rho^*\in L^1(\pi)$, it follows $\log \rho^* \in L^1(\pi)$.
\end{proof}

\end{document}
